\documentclass[11pt, oneside]{article}   	% use "amsart" instead of "article" for AMSLaTeX format
\usepackage{geometry}                		% See geometry.pdf to learn the layout options. There are lots.
\geometry{letterpaper}                   		% ... or a4paper or a5paper or ... 
%\geometry{landscape}                		% Activate for rotated page geometry
%\usepackage[parfill]{parskip}    		% Activate to begin paragraphs with an empty line rather than an indent
\usepackage{graphicx}				% Use pdf, png, jpg, or eps§ with pdflatex; use eps in DVI mode
								% TeX will automatically convert eps to pdf in pdflatex		
\usepackage{amssymb}
\usepackage{listings}
\usepackage{xcolor}
\usepackage{hyperref}

% Code listing style
\lstset{
    basicstyle=\ttfamily\small,
    breaklines=true,
    frame=single,
    backgroundcolor=\color{gray!10},
    keywordstyle=\color{blue},
    commentstyle=\color{green!60!black},
    stringstyle=\color{red},
    numbers=left,
    numberstyle=\tiny\color{gray}
}

\title{Weightlifting Coach \& Athlete Management Platform}
\author{Addy Cruz \\ Advisor: Dr. Helen Hu \\ CMPT 385 Senior Capstone}
\date{\today}

\begin{document}
\maketitle

\begin{abstract}
\noindent
This project proposes the development of a web based platform designed specifically for Olympic weightlifting coaches and athletes. The system will centralize training program management, athlete performance tracking, and competition day support into a single application. Unlike generic fitness apps, this platform is tailored to the unique workflows of weightlifting, addressing existing gaps in programming, meet management, and remote coaching. The platform combines software engineering, database design, and data science to provide a comprehensive solution that supports both coaches and athletes throughout the training and competition cycle. The technical implementation utilizes a modern full stack architecture with React frontend, Django REST Framework backend, PostgreSQL database, and Python based data science layer for predictive analytics and performance visualization.
\end{abstract}

\section*{Introduction}

This project addresses a significant gap in the current fitness technology landscape: the lack of domain specific tools for Olympic weightlifting. While numerous generic fitness applications exist, none provide the specialized features required for effective weightlifting coaching and athlete management. The proposed platform will serve as a comprehensive solution that bridges the gap between training program design, athlete performance tracking, and competition day operations.

The primary goal is to create an integrated system where coaches can manage multiple athletes simultaneously, storing and sharing daily, weekly, or monthly training programs while monitoring athlete adherence and progress. For athletes, the system will serve as a personal logbook, allowing them to record completed workouts, personal records, and feedback after each training session. This bidirectional communication enhances the coach athlete relationship and provides valuable data for performance analysis.

Beyond basic program management, the platform will generate graphical displays of progress, including training volume, intensity, personal record history, and Sinclair scores, a weightlifting specific metric used to compare lifters across different bodyweight classes. These visualizations will help both athletes and coaches understand readiness levels and long term strength development patterns.

A third major component focuses on competition day support, providing coaches with a real time tool for managing warm ups, attempt selections, and timing. This feature is particularly valuable for remote coaching scenarios, which have become increasingly common. The platform will also explore integration with existing meet software through a universal plugin or API, potentially serving as a proof of concept for standardized competition data management guidelines.

\section*{Background}

Olympic weightlifting is a sport that requires precise programming, detailed performance tracking, and strategic competition management. Coaches must balance multiple athletes' training cycles, monitor progress through various metrics, and make critical decisions during competitions. Current solutions in the market are fragmented: generic fitness apps lack weightlifting specific features, while specialized tools often focus on single aspects (programming OR tracking OR competition management) rather than providing an integrated solution.

The weightlifting community has expressed a need for tools that:
\begin{itemize}
    \item Support the unique structure of weightlifting programs (snatch, clean \& jerk, accessory work)
    \item Track performance metrics specific to the sport (Sinclair scores, competition totals, training volume)
    \item Provide competition day support for coaches managing multiple athletes
    \item Enable remote coaching capabilities
    \item Generate meaningful analytics from training data
\end{itemize}

This project draws from over a decade of personal experience in competitive weightlifting, including appearances at the American Open, bronze medal at the Praxis Cup, and current ranking in the top 4 of the Men's 85kg Master's 35 to 40 Division. This domain expertise informs the design decisions and ensures the platform addresses real world needs of the weightlifting community.

\section*{Proposed Work}

The system will follow a full stack development model, combining modern web technologies with data science capabilities to create a comprehensive platform. The architecture is designed for scalability, maintainability, and cost effectiveness, utilizing free tier cloud services suitable for a capstone project.

\subsection*{Technical Stack}

\textbf{Frontend:}
\begin{itemize}
    \item \textbf{React} with Vite for fast development and optimized production builds
    \item \textbf{Chart.js} and \textbf{react chartjs 2} for data visualization
    \item \textbf{React Router} for navigation and routing
    \item Modern UI components with responsive design
\end{itemize}

\textbf{Backend:}
\begin{itemize}
    \item \textbf{Django REST Framework} for robust API development
    \item \textbf{SimpleJWT} for token based authentication
    \item Role based access control (coach vs. athlete permissions)
    \item RESTful API design for clean separation of concerns
\end{itemize}

\textbf{Database:}
\begin{itemize}
    \item \textbf{PostgreSQL} for structured data storage
    \item Relational models for core entities (Users, Programs, Workouts)
    \item Event log design for historical training data
    \item Hybrid approach: relational core + event log for analytics
\end{itemize}

\textbf{Data Science Layer:}
\begin{itemize}
    \item \textbf{Python} with \textbf{pandas} for data analysis
    \item \textbf{scikit learn} for predictive modeling
    \item Sinclair score calculator (weightlifting specific metric)
    \item Performance analytics and visualization
    \item Future: ML models for attempt prediction and training load recommendations
\end{itemize}

\textbf{Deployment:}
\begin{itemize}
    \item \textbf{Frontend:} Netlify (free tier) with automatic GitHub deployments
    \item \textbf{Backend:} Render (free tier) for Django REST API
    \item \textbf{Database:} Supabase or Neon (free tier PostgreSQL)
    \item \textbf{Version Control:} GitHub with CI/CD pipelines
\end{itemize}

\subsection*{Specific Tasks}

\textbf{Core Functionality (MVP):}
\begin{enumerate}
    \item User authentication and role management (coach/athlete accounts)
    \item Training program creation and assignment system
    \item Athlete logbook for recording workouts and personal records
    \item Basic visualization dashboard showing progress over time
    \item Sinclair score calculator implementation
\end{enumerate}

\textbf{Advanced Features (Stretch Goals):}
\begin{enumerate}
    \item Competition day dashboard with timers and attempt management
    \item Real time warm up planning and timing calculations
    \item Integration with external meet software (proof of concept)
    \item Machine learning models for attempt selection recommendations
    \item Predictive analytics for peaking and tapering cycles
\end{enumerate}

\subsection*{Rationale}

The technology stack was selected based on several key criteria:

\textbf{React + Vite:} React is the most widely used frontend framework with extensive community support and documentation. Vite provides fast development and optimized production builds, making it ideal for rapid iteration during development.

\textbf{Django REST Framework:} Django is a mature, battle tested framework used by major companies (Instagram, Spotify, NASA). REST Framework provides excellent API development tools with built in authentication, serialization, and browsable API documentation. The framework handles complex requirements efficiently while maintaining code organization.

\textbf{PostgreSQL:} PostgreSQL is a robust, open source relational database that excels at handling complex queries and maintaining data integrity. Free tier options (Supabase, Neon) provide full PostgreSQL compatibility without vendor lock in, making it cost effective for a capstone project.

\textbf{Python Data Science Stack:} pandas and scikit learn are industry standard libraries for data analysis and machine learning, used by data scientists worldwide. Both libraries have extensive documentation and active communities, ensuring reliable implementation of analytics features.

\textbf{Free Tier Deployment:} All selected services offer free tiers specifically designed for development and small applications. This approach allows the project to demonstrate full stack deployment capabilities without incurring costs, while still providing production quality hosting.

\subsection*{Plan of Work}

\textbf{Phase 1: Foundation (Weeks 1 to 4)}
\begin{itemize}
    \item Set up development environment and project structure
    \item Implement user authentication with JWT tokens
    \item Create database models for core entities
    \item Build basic API endpoints for programs and workouts
    \item Develop initial React frontend with routing
\end{itemize}

\textbf{Phase 2: Core Features (Weeks 5 to 8)}
\begin{itemize}
    \item Implement training program creation and assignment
    \item Build athlete logbook interface
    \item Create workout logging functionality
    \item Develop basic dashboard with progress visualization
    \item Implement Sinclair score calculator
\end{itemize}

\textbf{Phase 3: Advanced Features (Weeks 9 to 12)}
\begin{itemize}
    \item Enhance visualization with Chart.js
    \item Implement performance analytics
    \item Develop competition dashboard (if time permits)
    \item Add predictive modeling features (if time permits)
    \item Conduct user testing and refinement
\end{itemize}

\textbf{Phase 4: Deployment and Documentation (Weeks 13 to 14)}
\begin{itemize}
    \item Deploy frontend to Netlify
    \item Deploy backend to Render
    \item Set up database on Supabase/Neon
    \item Configure CI/CD pipelines
    \item Complete documentation and final report
\end{itemize}

\section*{Timeline}

\begin{table}[h]
\centering
\begin{tabular}{|l|l|}
\hline
\textbf{Phase} & \textbf{Timeline} \\
\hline
Foundation Setup & Weeks 1 to 4 \\
Core Features Development & Weeks 5 to 8 \\
Advanced Features & Weeks 9 to 12 \\
Deployment \& Documentation & Weeks 13 to 14 \\
\hline
\end{tabular}
\caption{Project Development Timeline}
\end{table}

The project follows a 14 week development cycle, with the first 8 weeks focused on MVP completion and the remaining 6 weeks allocated for advanced features, testing, and deployment. This timeline allows for iterative development and refinement based on testing feedback.

\section*{Expected Deliverables}

\textbf{Minimum Viable Product (MVP):}
\begin{itemize}
    \item Fully functional web application with coach and athlete accounts
    \item Training program creation and assignment system
    \item Personal athlete logbook for recording workouts
    \item Basic visualization dashboard showing progress metrics
    \item Deployed application accessible via public URL
\end{itemize}

\textbf{Stretch Deliverables:}
\begin{itemize}
    \item Competition day dashboard with timers and attempt management
    \item Proof of concept for integrating external meet software
    \item Machine learning models for attempt selection recommendations
    \item Advanced analytics and predictive modeling features
\end{itemize}

\section*{Anticipated Impact}

This project will deliver a domain specific solution that supports Olympic weightlifting beyond what existing fitness platforms currently offer. It combines software engineering, database design, and data science into a single application while directly addressing challenges faced by both athletes and coaches. By uniting training management with competition support and analytics, the platform has the potential to make a meaningful impact on the weightlifting community.

For the student, this project provides an opportunity to demonstrate applied skills in building end to end systems while drawing from over a decade of personal experience in the sport. The technical challenges span multiple domains: full stack web development, database design, API development, data visualization, and machine learning, providing comprehensive experience relevant to modern software engineering careers.

\section*{Conclusion}

The Weightlifting Coach \& Athlete Management Platform represents a comprehensive solution to real world challenges in the weightlifting community. By leveraging modern web technologies, robust database design, and data science capabilities, the platform will provide coaches and athletes with tools that enhance training effectiveness and competition performance. The project demonstrates full stack development capabilities while addressing domain specific needs that existing solutions fail to meet.

The technical approach utilizes proven, industry standard technologies that integrate effectively, ensuring a reliable and maintainable system. The free tier deployment strategy makes the project cost effective while still demonstrating production quality deployment practices. Through this project, valuable experience is gained in software engineering, database design, API development, data visualization, and machine learning, skills highly relevant to modern technology careers.

\end{document}
